% ID: SEL-D0040-OB-E01
% ESTADO: validada. S041, cotejada con examen y criterios.
\textbf{a)} En el primer cuadrante, el recinto queda por
debajo de $y=4-x/3$ y por encima de $y=(x-8)/3$
y $y=(15-3x)/2$.
\begin{center}
\setlength{\unitlength}{8mm}
\begin{picture}(14,6)(-1,-1)
\qbezier(0,0)(5.5,0)(11,0)\put(11.2,0){$x$}
\qbezier(0,0)(0,1.8)(0,3.6)\put(0,3.7){$y$}
\qbezier(3,3)(4,1.5)(5,0)
\qbezier(3,3)(6.5,2.5)(10,2)
\qbezier(10,2)(9,1)(8,0)
\put(3,3.2){\makebox(0,0)[b]{$(3,3)$}}
\put(10.2,2){$(10,2)$}
\put(5,-0.3){\makebox(0,0)[t]{$(5,0)$}}
\put(8,-0.3){\makebox(0,0)[t]{$(8,0)$}}
\end{picture}
\end{center}
\textbf{b)} En $y=0$ se obtiene $5\le x\le8$.
La frontera superior corta a $3x+2y=15$ en $(3,3)$
y a $x-3y=8$ en $(10,2)$. Las dos fronteras inferiores
se cruzan por debajo del eje horizontal. Los vértices son
$(5,0)$, $(8,0)$, $(10,2)$ y $(3,3)$.

\textbf{c)} Al evaluar la función lineal en los vértices,
se obtienen $25$, $40$, $68$ y $42$, respectivamente.
El máximo es $\boxed{68}$ y se alcanza sólo en $(10,2)$.
